\documentclass[11pt,a4paper]{article}
\usepackage[margin=25mm]{geometry}
\usepackage{amsmath,amssymb,amsthm,booktabs}
\newtheorem{theorem}{Theorem}
\newcommand{\Sz}{\operatorname{Sz}}
\newcommand{\Tr}{\operatorname{Tr}}
\newcommand{\F}{\mathbb F}
\title{An order-five overlap with a Suzuki face over eight elements}
\author{Human Compression project: bounded F003 continuation}
\date{}
\begin{document}
\maketitle

This note treats the odd-order case complementary to the preceding
exclusion for a face $\Sz(8)$. We retain the entire preceding note,
including its coordinate conventions, characteristic-polynomial lemma,
linear equations, and exact certificate specification. Its group-theoretic
predecessors and classical interfaces remain retained as well.
The new work is a complete orbit reduction for an order-five middle
rotation and its polynomial certificates. The required field condition
is essential: no result is asserted here when the degree of the other
face is divisible by three.

\begin{theorem}\label{main}
Let $S=\Sz(2^N)$, with $N$ odd, and suppose
$a,b,c\in S$ satisfy the following conditions:
\begin{enumerate}
\item $a,c$ have odd order greater than one, $b$ has order five,
and $ab,bc,abc$ are involutions;
\item $\langle a,b\rangle$ is a Suzuki subfield group isomorphic
to $\Sz(8)$, and $\langle b,c\rangle$ is a Suzuki subfield group
isomorphic to $\Sz(2^e)$, where $\gcd(3,e)=1$.
\end{enumerate}
Then no such configuration exists.
\end{theorem}

Neither an intersection condition nor generation of $S$ is needed
for this assertion. Together with the accepted cyclic-overlap result
and the preceding order-four exclusion, it excludes both possible
middle-rotation orders for two genuine faces with cyclic intersection,
one face $\Sz(8)$, and coprime face degrees. This remains a partial
recovery of the fixed F003 source.

\section*{The first face and its orbits}

Conjugate the first face to its standard embedding over $K=\F_8$.
Use the matrices $U(x,y)$ and $J$ of the preceding note, and put
\[
                         b_0=z(1)J.
\]
This matrix has entries in $\F_2$ and order five, by direct
multiplication. Every element of order five in $\Sz(8)$ is conjugate
to it. Indeed $|\Sz(8)|=8^2(8^2+1)(8-1)$ has precisely one factor
five, so its subgroups of order five are conjugate by Sylow's theorem.
Within the standard $\Sz(2)=C_5:C_4$, the complement of order four
acts faithfully, and hence transitively, on the four generators of
$\langle b_0\rangle$. In the displayed representation,
$U(1,0)^{-1}b_0U(1,0)=b_0^2$, which verifies this transitivity
directly. We may therefore take
$b=b_0$. Put $u=ab$ and $v=bc$.

There are exactly 455 involutions in this standard $\Sz(8)$. They
are the seven matrices $z(r)$, $r\in K^\times$, and the matrices
\[
 U(x,y)^{-1}Jz(r)JU(x,y),\qquad x,y\in K,\quad r\in K^\times.
                                                               \tag{1}
\]
The first seven fix infinity. The others are grouped by their
distinct fixed points, and their nonzero centre parameters are unique
at each fixed point. The regular root-group action off infinity
therefore gives precisely $7+8^2\cdot7=455$ involutions, without
any omitted conjugacy class.

Since $bu=a^{-1}$ has odd order greater than one, the retained
characteristic-polynomial lemma implies $\Tr(bu)\ne0$.
Exactly 390 matrices in (1) and the first seven satisfy this
inequality. Conjugation by $b$ and entrywise squaring both preserve
the inequality and the equation $b=b_0$. They commute, and generate
an action of $C_5\times C_3$ on these 390 matrices. All its orbits
have size fifteen, giving twenty-six cases. This is a reduction by
explicit automorphisms; no assertion that these automorphisms give
the entire centralizer action is needed.

The exact checker reconstructs all 455 matrices from (1), verifies
their distinctness, selects all 390 nonzero traces, and closes every
orbit under both operations. The certificate file for each case
contains its representative and the full orbit. Their disjoint union
is verified to be the complete set of 390 matrices. Thus this
calculation covers every possible $u$, even if some listed pair were
to generate a smaller group. Such pairs are not discarded.

\section*{Descent and the polynomial identities}

Put $L=\F_{2^e}$ in the ambient field. The hypothesis on $e$ gives
\[
               K\cap L=\F_2,\qquad [KL:L]=3.          \tag{2}
\]
As before, set
\[
 B=\Tr(b^2v),\quad Q=\Tr(bv),\quad
 B_1=B^\theta,\quad Q_1=Q^\theta.
\]
All four lie in $L$, because $b,v$ belong to the second face and
traces are invariant under simultaneous conjugation. The equality
$Q=\Tr(c)$ and the odd order of $c$ give $Q\ne0$, hence $Q_1\ne0$.
We impose no nonvanishing condition on $B$ or $B_1$.

Write a general symplectic involution in the self-adjoint shape
\[
 V=\begin{pmatrix}
 v_0&v_1&v_2&v_3\\v_4&v_5&v_6&v_2\\
 v_7&v_8&v_5&v_1\\v_9&v_7&v_4&v_0
 \end{pmatrix}.
\]
For each of the twenty-six representatives $u$, impose
\[
 ub^{-1}V=Vbu,\qquad \Tr(b^2V)=B,\qquad \Tr(bV)=Q.    \tag{3}
\]
The coefficient matrix in the ten entries of $V$ has rank eight.
As in the predecessor, every solution is an explicit matrix
$V(B,Q,s_0,t_0)$, linear over $K$ in two free coordinates and the
two traces. Each coefficient matrix, its rank and the complete
substitution are checked by exact arithmetic.

Take the nonzero entries of $V^2-1$, followed by
$c_2(b^2V)-B_1$ and $c_2(bV)-Q_1$, as generators $I_i$.
They vanish for a hypothetical configuration. The retained
characteristic-polynomial lemma justifies the last two generators
for an order-five $b$ as well as for an order-four $b$.
The files supply every generator and exact identities
\[
             E_j=\sum_i C_{ij}I_i,
           \qquad E_j\in K[B,Q,B_1,Q_1].              \tag{4}
\]
Write each $E_j$ on the basis $1,z,z^2$, where $z^3+z+1=0$.
Its three coefficients are polynomials over $\F_2$ in the four
trace variables. Their values lie in $L$, so (2) implies that all
three vanish separately. Let $G_1,\ldots,G_{n-1}$ be the distinct
nonzero coefficient polynomials, and append
\[
                         G_n=wQQ_1-1.                \tag{5}
\]
The choice $w=(QQ_1)^{-1}\in L$ makes this last generator vanish.
This differs from the saturation in the order-four note; the present
argument does not assume $B\ne0$.

In twenty-four cases the exported identity is
$1=\sum_i H_iG_i$. In the other two, numbered 17 and 21 with
indices starting at zero, it is
\[
                  Q_1^3+Q_1^2+1=\sum_i H_iG_i.       \tag{6}
\]
The first alternative is immediately impossible. For the second,
the retained explicit identities
\[
 (X^5+X^4+X^3+X)(X^3+X^2+1)=X^8+X,
 \qquad (X^3+X^2+1)+X(X^2+X)=1
\]
show that a root in $L$ belongs to $K\cap L=\F_2$ but cannot
belong to $\F_2$, also impossible. This proves Theorem~\ref{main}.

All twenty-six cases carry the full multipliers in (4) and (6),
or the corresponding unit identity. A checker using the same exact
sparse-polynomial arithmetic as the preceding note verifies every
identity, the rank computations and the full orbit coverage. Its
format and variable conventions are unchanged, except for (5).
The present certificate data and checker are separately bound; the
previous packet is immutable. No Gr\"obner-basis assertion or finite
ambient-group enumeration is a proof premise.

The new exclusion is uniform in $N$ and in odd $e$ coprime to three.
When $3\mid e$, (2) fails, and this coefficient argument gives no
exclusion. The cases with larger first-face fields also remain open
in this recovery. No whole-original, conjecture, Lean, or compression
credit is assigned. Distinct review of this author packet is pending.
\end{document}
